\documentclass[10pt]{article}
\usepackage[a4paper,margin=22mm]{geometry}
\usepackage[T1]{fontenc}
\usepackage{lmodern}
\usepackage[hidelinks]{hyperref}
\setlength{\parindent}{0pt}
\setlength{\parskip}{5pt}
\setlength{\emergencystretch}{2em}
\pagestyle{empty}

\begin{document}
\raggedright

\begin{flushright}
Department of Computer Science and Engineering\\
SAGE University\\
Indore 452020, Madhya Pradesh, India
\end{flushright}

The Editors\\
\emph{Journal of Engineering and Applied Science}\\
Springer Nature

\medskip
Dear Editors,

We submit the manuscript, ``Reliable comparison of severe felt-intensity
classifiers with versioned and dependent event data: a USGS DYFI benchmark,''
for consideration as a Research Article.

Engineering model selection can be invalidated by catalogue revision,
outcome-derived inputs, dependent events crossing data splits, temporal
contamination, or unresolved point-score differences. The manuscript asks
whether performance and model-superiority claims remain supportable when these
risks are controlled. It contributes a released qualification workflow:
candidates first pass schema, proxy, role, and temporal audits, then are scored
on a fixed holdout and compared with paired sequence-cluster uncertainty.

The frozen \(M\geq5\) snapshot contains 7,520 source records, 2,362 eligible
events, and a 430-event temporal holdout. The worked qualification rejects one
temporal-proxy candidate, supports all four admissible learned baselines over
no-skill, and leaves the two lowest-Brier models unresolved. It therefore shows
when a candidate is admissible and when an apparent advantage is insufficient
to name a winner. The random-versus-sequence contrast is small and reversed and
is reported descriptively.

The article fits the journal's explicit scope in artificial and machine
intelligence, computer science, computational methods, modelling and simulation,
seismic evaluation, and engineering research methods. It follows the
decision-oriented form of benchmark and engineering machine-learning comparison
articles already published by the journal.

Data, code, frozen results, and reproduction instructions are available at
\url{https://github.com/akssha74/dyfi-benchmark-reliability}. A clean clone
passes the released validators and benchmark test suite and reconstructs the
reported results. The scientific results remain those of immutable release
v1.0.27; the JEAS manuscript adds only generated engineering-use assets derived
from the same frozen inputs.

The manuscript is original and is not under consideration elsewhere. Both
authors approved the manuscript and this submission. They declare no competing
interests and no specific funding. The study uses aggregate public-domain event
metadata, involves no human participants, and redistributes no respondent-level
identifying information. Generative-AI assistance is disclosed and all assisted
work was reviewed and verified by the authors. The public repository contains a
non-peer-reviewed manuscript copy and source, disclosed as a preprint without a
DOI; the software and derived data retain their MIT and CC0 terms.

\medskip
Sincerely,

Akshay Sharma (corresponding author) and Lalji Prasad\\
Department of Computer Science and Engineering\\
SAGE University, Indore, India\\
\href{mailto:akssha74@gmail.com}{akssha74@gmail.com}

\end{document}
